% !TEX program = xelatex
\documentclass[12pt]{article}
\usepackage{ctex}
\usepackage{amsmath}
\usepackage{graphicx}
\usepackage{float}
\usepackage{booktabs}
\usepackage{geometry}
\geometry{a4paper, margin=2cm}

\title{\Huge \textbf{实验一 小信号调谐放大器}}
\author{专业：【在此输入专业全称】 \quad 学号：【在此输入学号】 \quad 姓名：【在此输入姓名】}
\date{}

\begin{document}

\maketitle

\section*{一. 实验目的}
% 这里填写实验对应的具体目的，分点列出
\begin{enumerate}
    \item 【实验目的1】
    \item 【实验目的2】
    \item 【实验目的3】
\end{enumerate}

\section*{二. 实验内容概要}

\vspace{1em}\noindent{\large \textbf{1. 实验原理概述：}}
\begin{itemize}
    \item \textbf{单调谐放大器}：【简述单调谐放大器原理】
    \item \textbf{双调谐回路放大器}：【简述双调谐放大器原理】
    \item \textbf{放大器动态特性}：【简述动态范围相关原理】
\end{itemize}

\vspace{1em}\noindent{\large \textbf{2. 核心实验内容：}}
\begin{enumerate}
    \item[(1)] \textbf{【关键内容1】}：【详细说明操作方法】
    \item[(2)] \textbf{【关键内容2】}：【详细说明操作方法】
    \item[(3)] \textbf{【关键内容3】}：【详细说明操作方法】
\end{enumerate}

\section*{三. 具体实验}

\subsection*{任务一：单调谐放大器幅频特性测量}

% 这里填写具体的实验任务一流程简述
【此处填写实验操作说明和测试环境配置描述】

\subsubsection*{1. 测试结果}

\textbf{a. 以下为扫频法的测试结果：}

\begin{figure}[H]
\centering
% \includegraphics[width=0.8\textwidth]{【请添加扫频法截图路径】}
\caption{扫频法测得的单调谐放大器幅频特性}
\label{fig:sweep_freq_result}
\end{figure}

\textbf{b. 示波器输出结果：}
\begin{figure}[H]
\centering
% \includegraphics[width=0.8\textwidth]{【请添加示波器截图路径】}
\caption{单调谐示波器截图}
\label{fig:osc_result}
\end{figure}


\textbf{c. 以下为点测法的测试结果及单调放大器的幅频特性曲线：}

\begin{table}[H]
    \centering
    \makebox[\textwidth][c]{
        \begin{tabular}{c*{14}{c}}
        \toprule
        % 在此处填入输入频率
        输入频率 $f$ (MHz) & 【频率1】 & 【频率2】 & 【...】 \\
        \midrule
        % 在此处填入输出幅度
        输出幅度 $U_{o}$ (mV) & 【幅度1】 & 【幅度2】 & 【...】 \\
        \bottomrule
        \end{tabular}
    }
    \caption{单调谐放大器幅频特性测试结果}
\end{table}

\begin{figure}[H]
\centering
% \includegraphics[width=0.8\textwidth]{【在此添加特性曲线图片路径】}
\caption{单调谐放大器幅频特性曲线}
\label{fig:single_tuned_curve}
\end{figure}

\subsubsection*{2.结果分析}

\textbf{测试数据与通频带分析：}
\begin{itemize}
    \item \textbf{单调谐幅频特性：} 【在此处填写对最大值及中心频率偏移等现象的分析】
    \item \textbf{通频带计算：} 【在此处填写通频带的计算过程和结果分析】
\end{itemize}

\subsection*{任务二：观察集电极负载对单调谐放大器幅频特性的影响}

\subsubsection*{1. 测试结果}

\begin{figure}[H]
    \centering
    \begin{minipage}[b]{0.48\textwidth}
        \centering
        % \includegraphics[width=\textwidth, height=3.5cm]{【请在此处替换为左侧图片路径】} 
        \caption{断开负载时的特性曲线}
    \end{minipage}
    \hfill
    \begin{minipage}[b]{0.48\textwidth}
        \centering
        % \includegraphics[width=\textwidth, height=3.5cm]{【请在此处替换为右侧图片路径】} 
        \caption{接通负载时的特性曲线}
    \end{minipage}
    \vspace{0.5em}
\end{figure}

\subsubsection*{2. 结果分析}
\begin{itemize}
    \item \textbf{有载品质因数与通频带的关系：} 【分析负载对Q值和通频带的具体影响】
    \item \textbf{增益变化：} 【分析并联谐振阻抗及电压增益的对应变化】
\end{itemize}

\subsection*{任务三：双调谐回路谐振放大器幅频特性测量}

% 这里填写具体的实验任务三流程简述
【此处填写调整开关及示波器测试方法的详细说明】

\subsubsection*{1. 测试结果}

\textbf{a. 扫频法输出结果及收集到的数据如下：}

\begin{table}[H]
    \centering
    \begin{tabular}{c*{7}{c}}
    \toprule
    % 根据具体数据长度调整
    输入频率 $f$ (MHz) & 【频率1】 & 【频率2】 & 【...】 \\
    \midrule
    \textbf{输出幅度 $U_{o}$ (mV)} & 【幅度1】 & 【幅度2】 & 【...】 \\
    \bottomrule
    \end{tabular}
    \caption{双调谐放大器幅频特性测试结果}
    \label{tab:double_tuned_results}
\end{table}

\begin{itemize}
    \item[b.] \textbf{凹陷点频率测量}：【在此处填写凹陷点现象及对应读数】
    \item[c.] \textbf{幅频特性曲线绘制}：【说明矩形系数与单调谐对比情况】
    
    \begin{figure}[H]
    \centering
    % \includegraphics[width=0.8\textwidth]{【在此添加双调谐曲线图片路径】}
    \caption{双调谐放大器幅频特性曲线}
    \end{figure}
    
    \item[d.] \textbf{耦合电容对特性的影响}：【说明调整电容后的现象】
\end{itemize}

\begin{figure}[H]
    \centering
    \begin{minipage}[t]{0.32\textwidth}
        \centering
        % \includegraphics[width=\textwidth, height=3.5cm]{【请在此处替换为扫频曲线图片路径1】}
        \caption{耦合电容减小扫频曲线}
    \end{minipage}
    \hfill
    \begin{minipage}[t]{0.32\textwidth}
        \centering
        % \includegraphics[width=\textwidth, height=3.5cm]{【请在此处替换为扫频曲线图片路径2】} 
        \caption{耦合电容为某一值时扫频曲线}
    \end{minipage}
    \hfill
    \begin{minipage}[t]{0.32\textwidth}
        \centering
        % \includegraphics[width=\textwidth, height=3.5cm]{【请在此处替换为扫频曲线图片路径3】} 
        \caption{耦合电容增大扫频曲线}
    \end{minipage}
    \vspace{0.5em}
    \caption{不同耦合情况下的幅频特性扫频曲线对比}
\end{figure}

\subsubsection*{2. 结果分析}

\textbf{测试数据特征分析：}
\begin{itemize}
    \item \textbf{双调谐峰谷特征：} 【分析导致“凹凸交替”的过耦合等现象】
    \item \textbf{选择性较单调谐优良：} 【分析增益平坦度及频带边缘陡峭的表现】
\end{itemize}

\subsection*{任务四：放大器动态范围测量}

% 这里填写具体的实验任务四流程简述
【此处填写系统配置与电压幅度增加过程说明】

\subsubsection*{1. 测试结果}

\begin{table}[H]
\centering
\begin{tabular}{c*{10}{c}}
\toprule
放大器输入 (mV) & 【数据1】 & 【数据2】 & 【...】 \\
\midrule
\textbf{放大器输出 (V)} & 【数据1】 & 【数据2】 & 【...】 \\
\textbf{电压放大倍数} & 【数据1】 & 【数据2】 & 【...】 \\
\bottomrule
\end{tabular}
\caption{放大器动态范围测试结果}
\end{table}

 \begin{figure}[htbp]
 \centering
%  \includegraphics[width=0.8\textwidth]{【请在此放入动态范围曲线图片路径】}
 \caption{放大器动态范围曲线}
 \end{figure}

\subsubsection*{2. 结果分析}

\textbf{测试数据及失真阶段分析：}
\begin{itemize}
    \item \textbf{线性动态范围分析：} 【分析小信号线性区内电压成正比放大表现】
    \item \textbf{失真与非线性特征：} 【分析大信号激励下的倍数衰退与波形顶切现象】
\end{itemize}

\section*{四. 总结}

【在此处填写完整的实验感悟与总结，包括对于单双调谐原理在工程上区别的体会、实验现象及数据特征背后原因的探究，仪器的使用技巧与注意事项总结等】

\end{document}
